% **************************************************************************
%           Testarten + Pyramide
% **************************************************************************
\subsection{Grundlagen zu Codetesting}
\label{ch:SdF:Legacy:sec:TestPyramide}

% Code-Coverage
Wie zuvor bereits erwähnt, werden Tests dafür genutzt, das Code-Verhalten von unberührtem Code bei einer Änderung beizubehalten und somit die Code-Qualität garantieren zu können. Fundamental für die Qualitätssicherung ist das Konzept der Code-Coverage (dt. Code-Abdeckung). Sie \glqq misst wie viel Prozent des Codes durch Tests abgedeckt werden\grqq{} [\cite{IleneBurnstein:2003} nach \cite{Oppel:2025}] und hilft somit, ungetsteten Code aufindig zu machen. Eine hohe Code-Coverage führt zu einer hohen Softwarequalität und trägt damit zur Reduktion von Fehlern bei [\cite{IleneBurnstein:2003} nach \cite{Oppel:2025}]. % S. 13

% viel Legacy-Code = viel Testcode
Erfolgt eine hohe Code-Abdeckung im gewachsenen Legacy Code, so ist  proportional ebenso viel Test-Code zu bewältigen.
% Tests sind auch Code
El-Morabea et al. betont \glqq so wie wir schlechten Code haben können, können wir auch schlechte Tests haben. Schließlich sind Tests auch Code. \grqq{} [übersetzt nach \cite{ElMorabea:2021}]

Software kann zur besseren Testbarkeit und somit Erreichung einer höheren Testabdeckung modularisiert werden (s. \autoref{ch:SdF:Legacy:sec:Methoden}). Beides ist Part des \textit{Test Driven Development} (TDD). Es wird ein fehlschlagender Test geschrieben, dann wird dieser zum Erfolg gebracht und schließlich wird der Code refactored, in Form von zum Beispiel Modularisierungstechniken. Die neuen Module können erneut getestet werden \cite{ElMorabea:2021, Stantchev:2010}.

% Testarten und Pyramide(?)
Tests werden in verschiedene Testphasen unterschiedlichen Umfangs unterschieden. Neben dem Manuellen Testen gibt es automatisierte Test. \textit{Unit}-Tests überprüfen eine Einheit oder ein Modul. \textit{Integration}-Tests kümmern sich um die Integration, also das Zusammenspiel zwischen einzelnen Systemen, während \textit{System}-Tests ganze Systemkomplexe auf ihr Zusammenspiel überprüfen. Ein Beispiel dafür sind UI-Tests (User-Interface, dt. Oberfläche). Weitere Tests sind confiramtion-tests (dt. Abnahme- oder auch \textit{Akzeptanz}-Tests). \cite{ElMorabea:2021, Oppel:2025, Pi:2020}.

Bezüglich des Umfangs wird in der Lehre die \textit{Test-Pyramide} verwendet. Sie besagt, dass je simpler ein Test ist, desto mehr sollte es sie geben. Je komplexer und aufwändiger ein Test ist, desto weniger werden benötigt. In dem Zusammenhang wird auch nach sogenannten \textit{Happy-Paths} (dt. glückliche Pfade, wie das System funktionieren sollte) und \textit{corner-cases} (dt. Eck-Fälle, Situationen die vermeintlich seltener erscheinen und Fehler werfen) unterschieden. \cite{ElMorabea:2021}

Es werden weitderhin die Testmethoden \textit{Whitebox} und \textit{Blackbox} unterschieden. Whitebox-Test wie Unit- und Intergration-Tests sind für die Entwickler, welche ein profundes Wissen über das geteste System haben. Sie können hinter die Strukturen blicken und stellen einzelne Funktionen sicher. Blackbox-Tests wie Systemtests testen eher oberflächlich und geben an, ob das erwartete verhalten eines systems aufgetreten ist. Beispielsweise wird eine Schnittstelle auf ihre Antworten getestet. \cite{Pi:2020}
